\documentclass[11pt]{article}

\usepackage[margin=1in]{geometry}
\usepackage{amsmath}
\usepackage{booktabs}
\usepackage{array}
\usepackage{hyperref}

\title{Assignment 2 Report: Sequential Decision Agent for Retail Replenishment}
\author{Mustafa Yousuf}
\date{}

\begin{document}

\maketitle

\section*{Problem Framing}

This project trains and evaluates a sequential decision agent for single-SKU retail inventory replenishment. Each episode represents a 60-day selling season. On each day, the policy observes a compact inventory state and chooses an order quantity. The business objective is to improve operating profit while avoiding excessive stockouts, holding costs, and end-of-season markdown losses.

The Markov decision process is defined as follows. The state is a tuple of inventory bin, pipeline inventory bin, recent demand bin, day type, and season phase. The action space is discrete: \(\{0, 5, 10, 15, 20, 30\}\) units ordered per day. The reward approximates daily operating profit:

\[
\begin{aligned}
r_t ={}& \text{sales revenue} - \text{purchase cost} - \text{fixed order cost} \\
&- \text{holding cost} - \text{stockout penalty} - \text{terminal markdown loss}.
\end{aligned}
\]

The transition dynamics receive due supplier orders, sample stochastic demand, calculate sales and lost sales, add new orders to the delivery pipeline, and advance the day.

\section*{Policies Evaluated}

The project compares three policies:

\begin{itemize}
    \item Random ordering, used as a sanity-check baseline.
    \item A reorder-point heuristic, used as the practical business baseline.
    \item A tabular Q-learning agent trained on the normal-demand simulator.
\end{itemize}

Each policy is evaluated over 250 seeded episodes in five scenarios: normal demand, demand spike, demand slump, supplier delay, and end-season pressure.

\section*{Results}

The learned Q-learning policy improves over random ordering in most scenarios, but it does not outperform the reorder-point heuristic. The heuristic has higher mean profit in every evaluated scenario and is more stable across repeated episodes.

\begin{table}[htbp]
\centering
\begin{tabular}{lrrrr}
\toprule
Scenario & Reorder Profit & Q-Learning Profit & Profit Gap & Q Service \\
\midrule
Normal demand & 3976.35 & 3734.07 & 242.28 & 0.7422 \\
Demand spike & 3881.67 & 2990.31 & 891.36 & 0.6351 \\
Demand slump & 3754.98 & 3360.22 & 394.76 & 0.7604 \\
Supplier delay & 3039.37 & 2869.95 & 169.42 & 0.6577 \\
End-season pressure & 3965.08 & 3690.37 & 274.71 & 0.7387 \\
\bottomrule
\end{tabular}
\caption{Mean episode performance for the reorder-point heuristic and Q-learning policy.}
\end{table}

The key behavioral pattern is that Q-learning carries less inventory, which reduces holding and terminal markdown costs, but this creates more stockout exposure. In normal demand, Q-learning has lower holding cost than the heuristic, but it has more stockout days and lower total profit. The demand-spike scenario is the clearest failure case: Q-learning fulfills only about 63.5 percent of demand and loses about 891 profit per episode relative to the heuristic.

A heuristic sensitivity check also shows that nearby reorder-point settings can outperform the selected baseline. This strengthens the conclusion that the learned policy is not yet competitive with simple, interpretable business rules.

\section*{Failure Analysis and Governance}

The main production risks are reward hacking, unsafe stockout behavior, simulator overfitting, and supplier instability. If the stockout penalty is too low, the learned policy can protect simulated profit by under-ordering. Strong simulator performance would also not guarantee production readiness because real demand can shift due to promotions, seasonality, competitors, or supplier disruptions.

Any future learned policy should be evaluated with service-level guardrails, maximum order caps, stockout alerts, monitoring for abnormal inventory positions, and human approval for unusual order recommendations. Shadow testing should only be considered after a learned policy beats tuned heuristic baselines on profit, service level, stockout behavior, and edge-case robustness.

\section*{Recommendation}

The Q-learning policy should not be deployed. The retailer should keep the reorder-point heuristic as the operating baseline. The learned policy is useful as an experiment, but it is not ready for live automation because it underperforms a realistic heuristic and behaves poorly under demand spikes.

\end{document}
